% Supplementary display supporting Chapter 7.
\chapter{Supplementary Inference Diagnostics}
\label{app:inference-diagnostics}

Chapter~\ref{chap:inference-results} reports the exact rejection frequencies
at the 1\%, 5\%, and 10\% levels in
Table~\ref{tab:inference-null-summary}.  The additional figure below isolates
the principal 5\% comparison and displays its conditional Monte Carlo
uncertainty.  It is retained for visual completeness but is placed here because
the numerical table communicates the same evidence more compactly.

\section{Oracle-calibrated size at the 5\% level}
\label{sec:app-inference-size-figure}

\begin{figure}[!htbp]
    \centering
    \includegraphics[width=0.96\textwidth]{calibration_comparison.pdf}
    \caption[Oracle-calibrated null rejection probabilities]{Empirical null rejection probability of
    $\mathcal T_{n,70}$ at $n=100$, using the oracle weighted chi-square 5\%
    critical value.  Each point is based on 5,000 finite-sample replications;
    vertical bars are normal-approximation intervals equal to the estimate
    plus or minus 1.96 binomial Monte Carlo standard errors.  They condition on
    the realised critical values and do not include uncertainty from the
    50,000 oracle draws.  Since only oracle calibration was selected, the
    figure compares empirical size with the nominal line rather than comparing
    alternative feasible calibration procedures.}
    \label{fig:app-inference-calibration-comparison}
\end{figure}

All three point estimates lie modestly above 5\%, matching the heavier
finite-sample upper tails documented in
Figure~\ref{fig:inference-null-calibration}.  The figure should therefore be
read as a finite-sample size assessment of the oracle reference, not as
evidence that oracle calibration is exact or available in an empirical
application.
